\documentclass[a4paper]{article}
\usepackage[T1]{fontenc}
\usepackage[utf8]{inputenc}
\usepackage{newcent}
\usepackage{helvet}
\usepackage{graphicx}
\usepackage[pdftex, pdfborder={0 0 0}]{hyperref}
\frenchspacing

\include{bml}
\title{Système Monticelli's}
\author{Thomas et Romain Monticelli}
\begin{document}
\maketitle
\tableofcontents

\section{Introduction}

\section{Ouverture de 1m}

\subsection{Le Walsh}

À partir de 2 cartes. Le Walsh, basé sur le principe ``la majeure d'abord'', consistant à ne présenter les \d\ que dans les mains fortes ou sans majeures, est utilisé.
En cas de fit fort sans majeure, le fit mineur inversé est utilisé.

\begin{bidtable}
1\c---\\
1\d \> 3 possibilités :\\
\>4+\d\ et inv+\\
\>5+\d , 3-\h \s\ et 6-10H\\
\>BAL 5-7H et 3-\h \s \\
1\h \> \h /4+ (possibilité de \d\ plus longs que \h\ si <11h)\\
1\s \> \s /4+ (possibilité de \d\ plus longs que \s\ si <11h)\\
1NT \> 8-10 3-\h \s \\
2\c \> FMI : fit \c , 10+ (voir développements plus loin)\\
2\d \> 6+\d\ (et 2GH), 18H+\\
2\h \> TODO\\
2\s \> TODO\\
2NT	 \> 11-12 BAL (3\c\ NF, reste FM)\\
3\c	 \> 5\c , main irrégulière ~5-9 H\\
3X	 \> Barrage 7 cartes\\
3NT	 \> 12-14 BAL
\end{bidtable}

\subsection{La séquence 1\pdfc\ - 1\pdfd}

\begin{bidtable}
1\c-1\d;\\
1M \> 12-17 H, M/4 Main irrégulière (5\c )\\
1\h\+\\
1\s \> 4eme couleur forcing\\
\>(L’ouvreur réagit comme sur une enchère NAT mais ne peut pas dire 4\s ) \+\\
2\s \> 12-14\\
3\s \> 15+\+\\
3NT \> SO\\
4X \> cue fit \s \-\-\\
2\s \> 5\d 4\s\ FM\-\\
1NT \> 12-14 H, BAL (suite, TODO double 2)
\end{bidtable}

\subsection{La séquence 1\pdfc\ - 1M}

\begin{bidtable}
1\c-1M;\\
1NT \> suite TODO double 2\\
2\c\+\\
2\d \> 3eme couleur forcing\-
\end{bidtable}

\subsection{La séquence 1m - 1M}

\begin{bidtable}
1m-1M;\\
2NT \> pas de fit\+\\
3\c \> réponses dans l'ordre, style Baron\+\\
3\d \> 5m\\
3M \> 3M\\
\>3M' 4M'\\
\>3N aucun des cas précédents\-\\
3\d \> fit dans la m (FM ou TDC)\\
3M \> NAT NF\\
3M' \> BIC (55 ou 65) FM\\
3NT \> SO\\
4m/M' \> cue avec atout M\\
4\h \> (sur 1\s ) BIC NF\\
4M \> SO\\
4NT \> quantitatif\-\\
3NT \> fit main régulière\\
4(3)M’ \> splinter\\
4m’ \> splinter\\
4m \> 5+m 4+M\\
4M \> main extrêmement irrégulière avec peu de points d’honneurs
\end{bidtable}

\subsection{La séquence 1m - 2\pdfh}

\subsection{La séquence 1m - 2\pdfs}

\subsection{La séquence 1\pdfd\ - 2\pdfc}

\begin{bidtable}
1\d-2\c;\\
2\d \> irrégulier 11-15, BAL 12-13\\
2M \> inversée NAT (ou 4-4-4-1)\\
2NT \> BAL 14 ou 18-19\+\\
3\c \> relais\+\\
3\d \> 18-19, le reste NAT 14\-\-\\
3\c \> 16+ PJS Fit 3ème (plutôt irrég)\\
3\d \> 6+\d , 16+ problème dans une ou deux M\\
3M \> 65 sans inversée\\
3NT \> 6+\d , 16+ avec les arrêts M
\end{bidtable}

\subsection{Le fit mineur inversé}

\begin{bidtable}
1\c-2\c \> --\\
2NT \> 12-14- BAL\+\\
3\c \> SO\\
3\d \> arrêt \d \+\\
3\s \> arrêt \s \+\\
4\c \> SO, pas d'arrêt \h \-\-\\
Autre \> : arrêt\\
Sauts \> : courtes\+\\
4NT \> : to play (des points dans les courtes)\-\-\\
2\d \> 14+, arrêt \d , arrêts \h /\s\ possibles\+\\
2\h \> arrêt \h\ et pas mini\\
2\s \> arrêt \s , pas d'arrêt \h\ et pas mini\\
2NT \> arrêts des autres couleurs mais mini\\
3\c \> pas d'arrêts des autres couleurs et mini\\
3\d \> singleton (ou chicane) \d \-\\
2\h \> 14+, arrêt \h , pas d'arrêt \d , arrêt \s\ possible\\
2\s \> 14+, arrêt \s , pas d'arrêts \h /\d \\
3\c \> 12-14- irrégulier\+\\
Sauts \> : courtes\+\\
4NT \> : to play (des points dans les courtes)\-\-\\
3\d \> 14+, court \d , \c\ mieux que Vxxx\\
3\h \> 14+, court \h , \c\ mieux que Vxxx\\
3\s \> 14+, court \s , \c\ mieux que Vxxx\\
3NT \> 18-19 Bal\\
4\d \> 15-18, chicane \d , beaux \c \\
4\h \> 15-18, chicane \h , beaux \c \\
4\s \> 15-18, chicane \s , beaux \c 
\end{bidtable}

\begin{bidtable}
1\d-2\d \> --\\
2NT \> 12-14- BAL\+\\
3\c \> arrêt \c \+\\
3\s \> arrêt \s \+\\
4\d \> SO, pas d'arrêt \h \-\-\\
3\d \> SO\\
Autre \> : arrêt\\
Sauts \> : courtes\+\\
4NT \> : to play (des points dans les courtes)\-\-\\
2\h \> 14+, arrêt \h , arrêts \s /\c\ possibles\+\\
2\s \> arrêt \s\ et pas mini\\
2NT \> arrêts des autres couleurs mais mini\\
3\c \> arrêt \c , pas d'arrêt \s\ et pas mini\\
3\d \> pas d'arrêts des autres couleurs et mini\\
3\h \> singleton (ou chicane) \h \-\\
2\s \> 14+, arrêt \s , pas d'arrêt \h , arrêt \c\ possible\\
3\c \> 14+, arrêt \c , pas d'arrêts \h /\s \\
3\d \> 12-14- irrégulier\+\\
Sauts \> : courtes\+\\
4NT \> : to play (des points dans les courtes)\-\-\\
3\h \> 14+, court \h , \d\ mieux que Vxxx\\
3\s \> 14+, court \s , \d\ mieux que Vxxx\\
3NT \> 18-19 Bal\\
4\c \> 14+, court \c , \d\ mieux que Vxxx\\
4\h \> 15-18, chicane \h , beaux \d \\
4\s \> 15-18, chicane \s , beaux \d \\
5\c \> 15-18, chicane \c , beaux \d 
\end{bidtable}

\subsubsection{Soutien au niveau de 3}

\begin{bidtable}
1m\+\\
3m \> NF, on peut jouer ce contrat\+\\
4m \> prolongation de barrage\\
3x \> demande d'arrêt dans cette couleur\\
4x \> courte\-\-
\end{bidtable}

\subsubsection{Après interventions}

Le FMI ne s'applique (passer par le surcontre)

\begin{bidtable}
1\d-X\+\\
2\d \> 6-9\\
2NT \> 10-11 fitté (Truscott)\\
3\d \> barrage 6+\-
\end{bidtable}

\begin{bidtable}
1\c-1\s\+\\
2\c \> fit 6-10\\
2\s \> 11+, pas 4\h , souvent fitté \c \\
2NT \> naturel 10-12 avec arrêt\\
3\c \> barrage\\
4\c \> barrage\-
\end{bidtable}

\subsection{Enchère du répondant après une inversée}

\section{Ouverture de 1M}

\subsection{Fit limite TODO}

\subsection{2NT Fitté par 3 cartes FM TODO finir}

\begin{bidtable}
1M-2NT\\
3\c \> 12-14 avec une courte\+\\
3\d \> Relais\+\\
3\h \> Court \c \\
3\s \> Court \d \\
3NT \> Court M'\-\\
3M \> TODO\-\\
3\d \> 12-14 Pas de courte\+\\
3M \> TODO\-
\end{bidtable}

\subsection{Le gazzili}

\begin{bidtable}
1\h-1\s--\\
1NT \> BAL 12-15 (2\c\ Roudi, Reste Nat)\\
2\c \> Gazilli : Soit NAT, soit Fort\+\\
2\d \> 8+PH, FM si ouvreur fort. Tout le reste est plus faible, sauf 3M.\+\\
2\h \> 5\h 4\c\ 12-16 NF. Ensuite, Passe ou NAT Inv+\\
2\s \> 3 cartes à \s \+\\
2NT \> Relais\+\\
3m \> 5\h 4m3\s \\
3\h \> 6\h 3\s\ 6 levées\\
3\s \> BAL 15-17 avec 3\s \\
3NT \> BAL 18-19 avec 3\s \-\-\\
2NT \> BAL (54m22 / 5332) 15-16 sans 3\s \\
3m \> 5\h 4m 17+ (ou 5\h 5m 5 perdantes)\\
3\h \> 6\h 3\s\ 6 levées\\
3\s \> 5+\h 4\s\ 15-17\+\\
3NT \> demande de courte\\
4x \> cue\-\\
3NT \> BAL (54m22 / 5332) 16+-17 sans 3\s \-\\
2\h \> Préférence\\
2\s \> 5\s\ court \h\ ou 6+\s \\
2NT \> 5\d 4\s\ court \h\ (ensuite 3m préférence NF).\\
3\c \> Préférence\\
3\d \> 6+\d \\
3\h \> Fit différé\\
3\s \> INV\-\\
2x \> NAT (Sur la répétition à 2\h , 3\h\ est Inv)\\
2NT \> 6\h /4x ou 6\h /3\s \+\\
3\c \> Relais\+\\
3\d \> 6\h 4\d \\
3\h \> 6\h 3\s \\
3\s \> 6\h 4\s \\
3NT \> 6\h 4\c \-\-\\
3m \> 55 4 perdantes\\
3\h \> 6+\h\ sans 3\s , 7 levées\\
3\s \> 5\h 4\s\ 18-20\+\\
3NT \> demande de courte\\
4x \> cue\-\\
3NT \> BAL 18-19 sans 3\s \\
4m \> Fit 20+PJS, Splinter\\
4\h \> Sign-off\\
4\s \> Fit distributionnel
\end{bidtable}

\begin{bidtable}
1\h-1NT--\\
2\c \> Gazilli : Nat ou 17+\+\\
2\d \> 8+ PH, FM si ouvreur fort. Reste plus faible\+\\
2\h \> 5\h 4\c\ 12-16\\
2\s \> 5\h 4m irrégulier 17+\+\\
2NT \> Relais\+\\
3\c \> 5\h 4\c \\
3\d \> 5\h 4\d \\
3\h \> 5\h 4\d 4\c \-\-\\
2NT \> Bal (54m22 / 5332) 16+-17\\
3m \> 55 5 perdantes\\
3\h \> 6+\h\ 6 levées\\
3\s \> 6\h 5\s\ 5 perdantes\\
3NT \> Bal (54m22 / 5332) 16+-17\-\\
2\h \> Préférence\\
2\s \> 5\c\ 4\d \\
2NT \> 5\d , 4/5\c \\
3\c \> Préférence\\
3\d \> 6+\d \-\\
2\d\h \> Nat 12-16\\
2\s \> Inversée (2NT modérateur)\\
2NT \> 6\h 4m 17+\+\\
3\c \> Relais\+\\
3\d \> 6\h 4\d \\
3\h \> 6\h 4\c \-\-\\
3m \> 55 4 perdantes\\
3\h \> 6+\h , 7 levées\\
3\s \> 65 4 perdantes\\
3NT \> Bal 18-19\\
4m \> 65 sans TDC\\
4\h \> Sign-off
\end{bidtable}

\begin{bidtable}
1\s-1NT--
\end{bidtable}

\begin{bidtable}
2\c \> Gazilli : Naturel ou Fort\+\\
2\d \> 8+ PH, FM si ouvreur fort. Reste plus faible\+\\
2\h \> 3+\h\ (4\h\ 17+, inversée avec 3\h , 6\s 3\h\ 6 levées ou Bal 15+ avec 3\h )\+\\
2\s \> Relais\+\\
2NT \> Bal 15-17 avec 3\h \\
3\c\d \> 5\s 4m3\h \\
3\h \> 5\s 4\h \\
3\s \> 6\s 3\h \\
3NT \> Bal 18-19 avec 3\h \-\-\\
2\s \> 5\s 4\c\ 12-14 NF\\
2NT \> Bal 15-16 sans 3\h \\
3m \> 54 17+ sans 3\h\ ou 55 5 perdantes\\
3\h \> 55 5 perdantes\\
3\s \> 6+\s\ sans 3\h , 6 levées\\
3NT \> Bal 16+-17 sans 3\h \-\\
2\h \> 5\h\ court \s\ ou 6+\h \\
2\s \> Préférence\\
2NT \> 5\d 4\h\ court \s\ (ensuite 3m préférence NF)\\
3\c \> Préférence\\
3\d \> 6+\d \-\\
2\d\h\s \> Nat\\
2NT \> 6\s 4x / 6\s 3\h\ 17+\+\\
3\c \> Relais\+\\
3\d \> 6\s 4\d \\
3\h \> 6\s 4\h \\
3\s \> 6\s 3\h \\
3NT \> 6\s 4\c \-\-\\
3\c\d\h \> 55 4 perdantes\\
3\s \> 6+\s\ sans 3\h , 17+\\
3NT \> Bal 18-19 sans 3\h \\
4x \> 65 sans TDC\\
4\s \> Sign-off
\end{bidtable}

Après une réponse faible sur un Gazilli, 3x est 5M4x 17+ ou 5M5x 5 perdantes (sauf 3m sur la réponse de 2NT et sur la réponse de 2\s\ après 1\h -1NT), 3M Inv, 2NT Inv avec réponse naturelle et non forcing, 4M clôture.

Sur le soutien simple à 2M, l'ouvreur utilise les enchères d'essai dans les courtes (2NT court \s\ sur 2\h ) et généralisé à 2\s\ (sur 2\h\ : nommer des forces, 2NT force \s ) ou 2NT (sur 2\s , nommer des forces), saut court et TDC, 2NT puis saut long et TDC, courte puis répétition chicane et TDC.

\section{Ouverture de 1NT}

\begin{bidtable}
1NT\+\\
2\c \> Stayman\\
2\d \> Texas \h \\
2\h \> Texas \s \\
2\s \> ambigu : 8-9 H REG ou Texas \c\ TODO developpements\\
2NT \> Texas \d \\
3X \> NAT 6 cartes, TDC\+\\
3X+1 \> intéressé, ambigu\-\\
3NT \> Pour jouer\\
4\c \> Bic mineur\+\\
4\d \> fit \d \\
4\h\s \> fit \c \\
4NT \> Coup de frein\-\\
4\d \> Bic M\\
4\h\s \> Pour jouer\\
4NT \> Quantitatif\\
5\c\d \> Pour jouer\-
\end{bidtable}

\subsection{Stayman}

Pas utilisé avec des main régulières, FM, 1 seule majeure 4e.
Utilisé avec une majeure 5e dans la zone 8-9H.

\begin{bidtable}
2\c \> TODO
\end{bidtable}

\subsection{1NT - 2\pdfd\ : Texas \pdfh\ TODO}

Remarques :

\begin{itemize}
\item Avec 4 cartes à \h\ et minimum, 4333 on rectifie le plus souvent à 3\h\ 

\item avec 4 cartes à \h\ et maximum, on rectifie à 2NT

\end{itemize}

\subsection{1NT - 2\pdfh\ : Texas \pdfs\ TODO}

Remarques :

\begin{itemize}
\item avec 4 cartes à \s\ et maximum, on rectifie à 2NT

\end{itemize}

\subsection{1NT-2C : Texas \pdfc\ TODO}

\subsection{1NT-2NT : Texas \pdfd\ TODO}

Remarque : L'ouvreur peut moduler son soutien, 3\c\ max fitté et 3\d\ RAS.
Il est préférable que l'ouvreur ait une tenue à \c\ lorsqu'il répond 3\c .

\section{Ouverture de 2\pdfc}

\subsection{Définition}

\begin{itemize}
\item FAIBLE en \d\ (constructif VUL et en 1ère et 2ème NV) et 11-13 en 4ème

\item 22-24 BAL

\item GAMBLING solide en \c\ 

\item FM sauf Bicolores Majeurs

\end{itemize}

\begin{bidtable}
2\c--\+\\
2\d \> relais\+\\
Pass \> FAIBLE \d \\
2\h \> 24+ balancé\\
\>Bicolore avec au moins 4\h \\
\>Unicolore \h\ balancé 22-25\+\\
2\s \> relais\+\\
2NT \> 25H et + (développements comme sur 2N)\\
3\c \> 5+\c , =4\h \\
3\d \> 5+\d , =4\h \\
3\h \> 5+\h , 4+\c \\
3\s \> 5+\h , 4+\d \\
3NT \> Unicolore \h\ Gambling (style régulier avec 6/7\h\ plein et 2 As)\\
4\c\d \> 6\h /5m moins fort que 3\h\ suivi de 4\c\ (5+/5)\\
\>Ex: \s\ 2 \h\ ADV987 \d\ 8 \c\ ARDT2\-\-\\
2\s \> Bicolore avec 4+\s \\
\>Unicolore \s\ balancé 22-25\+\\
2NT \> relais\+\\
3\c \> 5+\c , =4\s \\
3\d \> 5+\d , =4\s \\
3\h \> 5+\s , 4+\c \\
3\s \> 5+\s , 4+\d \\
3NT \> Unicolore \s\ Gambling\\
4\c\d \> 6\s /5m moins fort que 3\s\ suivi de 4\c\ (5+/5)\-\-\\
2NT \> 22+ - 24H\\
3\c\d \> Unicolores ou bicolore mineur (la collante est ambigue)\\
3\h\s \> Unicolores irréguliers\\
3NT \> Gambling \c \\
4\c\d \> Bicolores 6/5m \h\ (6\h -5\c /\d\ + faible que 2\c -2\d -2\h -2\s -4\c /\d )\\
4\h\s \> Bicolores 6/5m \s\ (pareil)\-\\
2\h\s3\c \> F1T (pas forcing en paires !)\\
\>Nouvelle couleur GF\+\\
2NT \> Max misfit\\
3\d \> Min\\
3x/4x \> Fit, Splinters mains faibles\-\\
2NT \> Relais F1 avec espoir de manche en face d’un faible à \d \+\\
3\c \> faible \d\ et une courte\+\\
3\d \> NF courte s’annonce par paliers si max\\
\>(=> Passe si min du faible, annonce la courte par palier si max du faible)\-\\
3\d \> faible \d\ sans courte\\
3\h \> Max pièce \c /\h , pas de courte\+\\
3P \> relais\+\\
3NT \> court \c \\
4\c \> court \h \-\-\\
3\s \> Max force sans courte\\
3NT \> beaux \d\ ARD ou ARVT sans courte\\
4X \> GF\-\\
3\d \> 7-11H (Forcing > 4NT si main forte, mêmes dvpts)\\
3\h\s4\c \> NAT + Fit\-
\end{bidtable}

\subsection{Avec les mains TRICOLORES, on reparle toujours sur 3NT.}

Ex : \s\ - \h\ ADV8 \d\ ARD98 \c\ ARV10
(2\c -2\d -2\h -2\s -3\d -3NT-4\c -...)

\subsection{Après interventions}

\begin{itemize}
\item Contre punitif

\item Toute enchère F1T

\item 2NT R fort

\item Cue demande d'arret

\end{itemize}

\subsection{Après contre}

\begin{itemize}
\item Passe des \c !

\item Surcontre FORT ou UNICOLORE

\item 2\d\ ambigu

\item Nouvelle couleur F1T

\end{itemize}

\section{Ouverture de 2\pdfd}

\subsection{Définition}

\begin{itemize}
\item Unicolore FI

\item Bicolore Majeur FI ou FM

\end{itemize}

\begin{bidtable}
2\d-2\h\\
2\s \> Montre un bicolore majeur \textbf{+ long à \s }\+\\
2NT \> Relais\+\\
3\h \> 5/5 MAJ FI 18-21 PH\\
3\s \> 6\s /4\h\ FI 18-21 PH\\
3NT \> 5422 FM\\
4\c\d \> Tricolore parfait\\
4\h \> 65 FM\\
4\s \> 74 FM\\
3\d \> 5431 FM\\
3\c \> FM ambigu\+\\
3\d \> Relais\+\\
3\h \> 55 FM\\
3\s \> 64 FM\\
3NT \> 5413 FM\-\-\-\\
3\c\d \> Aucun intérêt pour les majeures, naturel (pas belle 6e)\\
3\h\s \> TDC\\
4\h\s \> NF\-
\end{bidtable}

\begin{bidtable}
2\d-2\h\\
2NT \> Montre un bicolore majeur \textbf{+ long à \h }\+\\
3\c \> Relais\+\\
3\d \> Résidu mineur 31 ou 30\+\\
3\h \> Relais\+\\
3\s \> 4513\\
3NT \> 4531\\
4X \> 4603 ou 4630\-\-\\
3\h\\
3\s \> 65 FM\\
3NT \> 5422 FM\\
4\c\d \> Tricolore parfait\\
4\h \> 74 FM\-\-
\end{bidtable}

\begin{bidtable}
2\d-2\h\\
3\c\d\h\s \> ACOL naturel (=FI) Collante ambigue, éventuellement pour tenter de jouer 3NT de la main forte\+\\
3\c\+\\
3\s \> Ambigu\+\\
3NT \> pour jouer\+\\
4\h \> TDC avec cue \s \\
4NT \> Quantitatif\-\-\\
4X \> Cue, TDC fit \h \-\\
3NT \> 9-10 plis en \d\ avec au moins 2 grosses pièces extérieures\-
\end{bidtable}

\begin{bidtable}
2\d\+\\
2\s \> Naturel 5 cartes avec deux gros honneurs ou 5/5 correct\\
3\c\d\h\s \> Transferts pour unicolores Q+ avec 2 gros honneurs\\
\>Le soutien ou la rectification du texas montre un intérêt pour la couleur du répondant\\
2NT \> 5/5 Mineur FM\-
\end{bidtable}

\subsection{Après intervention en 2e}

\subsubsection{Contre}

\begin{itemize}
\item Punitif si intervention au palier de 2

\item Appel si intervention au palier de 3 ou plus

\end{itemize}

\subsubsection{Passe RAS}

Ensuite, contre de l'ouvreur est punitif au niveau de 2 et montre le bicolore majeur ou une demande d'arrêt au niveau de 3.

\subsubsection{3\pdfc\pdfd\pdfh\pdfs\ Transferts (Transfert cue - Tricolore court)}

\subsubsection{2NT Bicolore mineur}

\subsubsection{Après 2\pdfd\ - X}

\begin{itemize}
\item Passe : envie de jouer

\item Tout le reste sauf 2\h\ : Naturel

\item XX : Naturel (4 belles cartes, pas 5 ou 6)

\end{itemize}

\subsection{Après intervention en 4e}

\begin{itemize}
\item PASSE, CUE-BID et ENCHERES A SAUT montrent le BICOLORE MAJEUR JUSQU’À 3\s .

\end{itemize}

\begin{itemize}
\item CONTRE demande l’arrêt

\end{itemize}

\subsection{Corollaires (à vérifier avec Gazilli !)}

\begin{bidtable}
1\s-1NT\\
3\h \> montre un 5=/4= dans la zone 18-21H
\end{bidtable}

\begin{bidtable}
1\s-1NT\\
4\h \> montre un 6/5 de 14 à 17H
\end{bidtable}

\begin{bidtable}
1\h-1NT\\
3\s \> montre un 6/5 de 14 à 17H, c’est NF !
\end{bidtable}

\section{Ouverture de 2\pdfh\ et 2\pdfs}

Montre une main de 5 à 10PH comportant une couleur 6ème (qualité de la couleur dépend de la position et de la vulnérabilité).

\begin{bidtable}
2M\+\\
3/4M \> Prolongation de barrage, 4M est une main faible ou forte qui espère gagner la manche.\\
2/3X \> Forcing un tour, l'ouvreur donne le fit (xxx ou Hx), répète sa couleur ou dit 3SA.\\
2NT \> Relais fort\+\\
2/3X \> FM, Max avec gros honneur.\+\\
3M \> TDC\+\\
4Y \> singleton Y\\
4M \> pas de singleton\-\-\\
3M \> Mini\\
3NT \> ARDxxx\\
4x \> Belle main, singleton X\-\\
4X \> Rencontre ?\-
\end{bidtable}

Attention, ``La force prime sur le courte''

\section{Ouverture de 2NT}

Cette ouverture peut contenir une majeure 5e.

\begin{bidtable}
2NT---\\
3\c \> Puppet Stayman\\
3\d \> Texas \h \\
3\h \> Texas \s \\
3\s \> Texas pour 3N\\
3NT \> 5\s\ et 4\h \\
4\c \> Bicolore mineur TDC\\
4\d \> Bicolore majeur, limité à la manche ou certitude de chelem
\end{bidtable}

\subsection{Le Puppet Stayman}

\begin{bidtable}
2NT-3\c \> ---\\
3\d \> 4\h\ et/ou 4\s \+\\
3\h \> 4 cartes à \s \+\\
3\s \> fit \s \+\\
4X \> envie de chelem\-\\
3NT \> 4 cartes à \h\ (il est donc nécessaire de passer par 3\h\ avec\\
\>les deux majeures 4e et des envies de chelem)\\
\>4X fit avec envie de chelem\-\\
3\s \> 4 cartes à \h\ et pas 4 cartes à \s \\
3NT \> 3-\h\ et 3-\s \\
4\d \> 4\h\ et 4\s\ sans envie de chelem\-\\
3\h \> 5\h \\
3\s \> 5\s \\
3NT \> 3-\h\ et 3-\s 
\end{bidtable}

\subsection{Les Texas}

\begin{bidtable}
2NT---\\
3\d\+\\
3\h \> 2\h \\
3\s \> 5\s\ et 2\h \+\\
4\c/4\d \> ctrl \c /\d\ et fit \s \\
4\h/4\s \> Arrêt\-\\
3NT \> 3+\h\ et ctrl \s \\
4\c \> 3+\h\ et ctrl \c\ (sans ctrl \s )\\
4\d \> 3+\h\ et ctrl \d \\
4\h \> 3+\h\ et tous les ctrl\-
\end{bidtable}

\begin{bidtable}
3\h\+\\
3\s \> 2\s \\
3NT \> 5\h\ et 2\s \+\\
4\c/4\d \> ctrl \c /\d\ et fit \h \\
4\h/4\s \> Arrêt\-\\
4\c \> 3+\s\ et ctrl \c \\
4\d \> 3+\s\ et ctrl \d\ (sans ctrl \c )\\
4\h \> 3+\s\ et ctrl \h \\
4\s \> 3+\s\ et tous les ctrl\-
\end{bidtable}

\begin{bidtable}
3\s\+\\
3NT\+\\
4\c/4\d \> Envie ou certitude de chelem à \c /\d . L'ouvreur décourage par 4SA.\\
4\h/S \> Envie de chelem à \h /\s . L'ouvreur passe ou demande les clés.\\
\>Avec certitude de chelem, il faut passer par un texas.\-\-
\end{bidtable}

\subsection{Autres réponses}

\begin{bidtable}
4\c\+\\
4\d \> fit \d \\
4\h \> fit \c\ et ctrl \h \\
4\s \> fit \c , ctrl \s , pas de ctrl \h \-
\end{bidtable}

\begin{bidtable}
4\d \> L'ouvreur choisit sa majeure
\end{bidtable}

\begin{bidtable}
4NT \> Minimum, l'ouvreur passe.\\
\>Maximum, les réponses sont les suivantes :\+\\
5\c \> 1 ou 4 As\\
5\d \> 0 ou 3 As\\
5\h \> 2 As et 5C\\
5\s \> 2 As et 5D\\
5NT \> 2 As et pas de min 5e\-
\end{bidtable}

\section{Les barrages et ouvertures à haut palier}

\subsection{Namyats 4\pdfc/4\pdfd}

Cette ouverture montre une couleur \h /\s\ 8e fermée ou 7e fermée avec un as.
La rectification à la manche est un arrêt.

L'enchère juste au-dessus montre un intérêt pour le chelem et demande à l'ouvreur de:

\begin{itemize}
\item Nommer la manche sans as annexe

\item Nommer son as annexe

\end{itemize}

\subsection{3SA}

Mineure 8e fermée sans rien à côté. 
4/5\c \d\ PoC

\section{Enchères après passe}

TBD

\section{Enchères de chelem}

L'enchère de 4NT est un Blackwood à 5 clés incluant le roi d'atout. L'atout est la couleur explicitement ou implicitement (via cue-bid) fittée ou à défaut la dernière couleur nommée naturellement

\begin{bidtable}
4NT\+\\
5\c \> 4 ou 1 clé(s)\\
5\d \> 3 ou 0 clé(s)\\
5\h \> 2 clés sans dame d'atout\\
5\s \> 2 clés avec dame d'atout\\
5NT \> 1 clé et une chicane\\
6x \> 2 clés et chicane x utile (ou chicane au-dessus si x = atout)\-
\end{bidtable}

Sur une réponse 5m, le palier immédiatement supérieur en dehors de l'atout est un relais pour la demande de la dame d'atout et des rois spécifiques.
Avec la dame d'atout, on nomme économiquement la couleur du premier roi ou 5SA déniant un roi.

Le retour à l'atout au palier le plus économique dénie la dame d'atout. Ensuite, le palier au-dessus demande de nommer les rois dans un ordre économique.

Une autre enchère non-ambigüe explore pour le grand chelem, demandant un complément.

Blackwood d'exclusion : Double saut anormal avec fit précisé !

Une réponse de 5NT avec une zone initiale de 2 pts (par exemple l'ouverture de 2NT) propose de découvrir un fit 4-4 au palier de 6.

En cas de bicolore et sans fit précisé par manque de place, le Blackwood est à 6 clés :

\begin{itemize}
\item 5H : pas de dame

\item 5S : la dame la moins chère

\item 5NT : la plus chère

\item 6C : les deux

\end{itemize}

Après le BW, 5NT est une demande de roi :

\begin{bidtable}
5NT\+\\
Retour \> à l'atout : pas de roi.\\
6X \> : Roi X ou les deux autres.\\
6NT \> : Tous les rois.\-
\end{bidtable}

Sur un 4NT quantitatif :

\begin{bidtable}
4NT\+\\
Passe \> mini\\
5\c \> 4 ou 1 As\\
5\d \> 3 ou 0 As\\
5\h \> 2 As et 5\c \\
5\s \> 2 As et 5\d \-
\end{bidtable}

Intervention adverse sur le BW :

\begin{itemize}
\item Passe : 1er palier 

\item X ou XX : 2ème palier

\item Autre :  3ème palier et plus.

\end{itemize}

Intervention adverse sur la réponse de 5m au BW :

\begin{itemize}
\item Passe : demande de Dame d'atout

\item X ou XX : demande de nommer les Rois.

\end{itemize}

Intervention adverse par contre sur les cue-bids :

\begin{itemize}
\item Passe : sans contrôle

\item XX : contrôle du premier tour

\item Autre : contrôle du second tour (retour à l'atout = second tour sans autre cue).

\end{itemize}

Intervention adverse en couleur sur les cue-bids :

\begin{itemize}
\item X : punitif

\item Passe : sans contrôle

\item Autre : souvent court dans la couleur

\end{itemize}

Intervention adverse par contre sur un cue-bid dans une courte connue (Splinter par exemple) :

\begin{itemize}
\item Passe : pas de contrôle

\item XX : contrôle du premier tour

\item Autre : cue-bid + contrôle du second tour

\end{itemize}

\section{Défense TODO}

\subsection{Sur le SA fort}

\subsection{Sur le SA faible}

\section{Enchères compétitives}

\section{Jeu de la carte TODO
}

\end{document}
